\documentclass[article,nojss]{jss}

%% -- Packages --------------------------------------------------
\usepackage{amsmath,amssymb,amsfonts}
\usepackage{bm}
\usepackage{booktabs}
\usepackage{graphicx}
\usepackage{float}
\usepackage{xcolor}
\usepackage{enumitem}
\usepackage{hyperref}
\usepackage{algorithm}
\usepackage{algorithmic}

%% -- Page header -----------------------------------------------
\usepackage{fancyhdr}
\fancyhead[L]{}
\fancyhead[C]{}
\fancyhead[R]{\thepage}
\fancyfoot{}
\renewcommand{\headrulewidth}{0pt}

%% -- Metadata --------------------------------------------------
\title{Algorithm: ES-MDA with Adaptive Student-\emph{t} Likelihood}
\Plaintitle{ES-MDA with Adaptive Student-t Likelihood}
\Shorttitle{Student-t ES-MDA Algorithm}

\author{Frederick~Bennett}
\Plainauthor{Frederick Bennett}

\Abstract{
This document presents the complete algorithm for the adaptive
Student-$t$ ensemble smoother described in the accompanying paper.
The algorithm extends the standard ES-MDA framework with three
components: (1)~latent per-observation precision weights $\lambda_i$
drawn from a Gamma distribution, implementing the scale-mixture-of-normals
representation of the Student-$t$ likelihood; (2)~adaptive estimation of
the degrees of freedom $\nu$ from the ensemble residuals at each
iteration via kurtosis matching; and (3)~inflation of both the
observation perturbation by $1/\sqrt{\lambda_i}$ and of the noise
covariance matrix by $1/\lambda_i$, preserving the exactness of the
ensemble Kalman
update conditional on the latent weights.
The algorithm adds approximately 50 lines to an existing ES-MDA
implementation and introduces no new tuning parameters beyond those
already required by the base method.
}

\Keywords{ensemble smoother, Student-t, pseudocode, algorithm, ES-MDA}
\Plainkeywords{ensemble smoother, Student-t, pseudocode, algorithm, ES-MDA}

\Address{
  Frederick Bennett\\
  Department of Environment, Tourism, Science and Innovation\\
  Queensland Government\\
  E-mail: \email{frederick.bennett@des.qld.gov.au}
}

\begin{document}

%=============================================================================
\section{Algorithm Description}
%=============================================================================

The algorithm extends the standard ES-MDA loop with three insertions,
marked \textbf{(S1)--(S3)} in Algorithm~\ref{alg:esmda_student_t}.
The remaining steps are unchanged from the Gaussian ES-MDA formulation
of \cite{Emerick2013}.

\subsection*{(S1) Adaptive $\nu$ Estimation}

After the forward model has been evaluated for all ensemble members
(Step~2a), the ensemble-mean residuals are standardised by the current
noise estimate $\bar{\phi}_i$.
A robust sample kurtosis $\hat{\kappa}$ is computed from residuals
clipped to $\pm 8$ standard deviations, preventing individual extreme
outliers from dominating the fourth moment.
If $\hat{\kappa} > 3.1$ (indicating heavier-than-Gaussian tails), the
degrees of freedom are estimated via the method-of-moments formula
$\hat{\nu} = 4 + 6/(\hat{\kappa} - 3)$, which is exact for a
Student-$t$ distribution with $\nu > 4$.
Otherwise $\hat{\nu}$ is set to a high value (100), reflecting
approximately Gaussian behaviour.
The estimate is smoothed with an exponential moving average
($\eta = 0.7$) and clamped to $[3, 100]$.
At the first iteration, $\nu$ is initialised to a user-specified value
(default 8.0) since residuals from the prior ensemble are unreliable.

\subsection*{(S2) Latent Precision Weights}

Each observation $i$ receives a latent precision multiplier $\lambda_i$,
drawn from a Gamma distribution.
At iteration~1, the prior $\text{Gamma}(\nu/2, \nu/2)$ is used.
At subsequent iterations, the posterior conditional on the current
standardised residual $r_i$ is used:
$\lambda_i \mid r_i \sim \text{Gamma}((\nu+1)/2, (\nu+r_i^2)/2)$.
This posterior draw targets down-weighting at observations that currently
have large ensemble-mean residuals, rather than relying on random draws.
The parametrisation $\text{Gamma}(\nu/2, \nu/2)$ gives $\mathbb{E}[\lambda_i]
= 1$, so typical (well-fit) observations experience standard perturbation
variance.

\subsection*{(S3) Perturbation and Covariance Scaling}

The latent weights enter the ES-MDA machinery in two places, both
required for statistical consistency with the scale-mixture
representation.
First, the observation perturbation is scaled by $1/\sqrt{\lambda_i}$
(Step~2d), so that observations with $\lambda_i \ll 1$ receive
proportionally larger perturbation noise.
Second, the observation noise covariance matrix $C_d$ is inflated by
$1/\lambda_i$ (Step~2e), increasing the effective observation error
variance for outliers and shrinking their contribution to the Kalman
gain.
Both scalings use the same weight vector
$\mathbf{w} = 1/\sqrt{\max(\boldsymbol{\lambda}, \epsilon)}$ with
$\epsilon = 10^{-6}$ for numerical stability.

Conditional on $\boldsymbol{\lambda}$, the likelihood is exactly
Gaussian, so the Kalman algebra (covariance computation, subspace
inversion, parameter update) is unchanged.
The IKEA Laplace noise estimation (Step~2f) operates on raw residuals
and is unaffected by the latent weights, maintaining the separation
between baseline noise estimation ($\phi$) and outlier accommodation
($\boldsymbol{\lambda}$).

%=============================================================================
\section{Pseudocode}
%=============================================================================

\begin{algorithm}[tbp]
\caption{ES-MDA with Adaptive Student-$t$ Likelihood}
\label{alg:esmda_student_t}
\begin{algorithmic}[1]
\REQUIRE Prior ensemble $\mathbf{M}_0 \in \mathbb{R}^{N_m \times N_e}$,
         observations $\mathbf{y} \in \mathbb{R}^{N_d}$,
         inflation factors $\alpha_1, \ldots, \alpha_{N_a}$,
         initial $\nu_0 = 8$, smoothing $\eta = 0.7$
\ENSURE Posterior ensemble $\mathbf{M}_{N_a}$, final $\nu$, final $\boldsymbol{\phi}$

\STATE $\boldsymbol{\phi} \gets \text{HalfNormal}(\mathbf{y}, \text{error})$
       \hfill $\triangleright$ per-group noise prior
\STATE $\nu \gets \nu_0$
\STATE \textbf{for} $t = 1$ \textbf{to} $N_a$ \textbf{do}

\STATE \quad $\tilde{\mathbf{M}} \gets \text{inv\_transform}(\mathbf{M})$
       \hfill $\triangleright$ to physical parameter space
\STATE \quad $\mathbf{D} \gets \text{fill\_ensemble}(\tilde{\mathbf{M}})$
       \hfill $\triangleright$ $\mathbf{D} \in \mathbb{R}^{N_d \times N_e}$

\STATE \quad \textbf{(S1) Estimate $\nu$} \textbf{(if $t > 1$)}
\STATE \quad\quad $\bar{\mathbf{r}} \gets \mathbf{y} - \text{mean}(\mathbf{D}, \text{axis}=1)$
\STATE \quad\quad $\hat{\mathbf{r}} \gets \bar{\mathbf{r}} \;/\; \max(\text{mean}(\boldsymbol{\phi}, \text{axis}=1),\, 10^{-8})$
\STATE \quad\quad $\hat{\mathbf{r}} \gets \text{clip}(\hat{\mathbf{r}}, -8, 8)$
\STATE \quad\quad $\hat{\kappa} \gets \mathbb{E}[\hat{\mathbf{r}}^4] \;/\; \mathbb{E}[\hat{\mathbf{r}}^2]^2$
\STATE \quad\quad \textbf{if} $\hat{\kappa} > 3.1$ \textbf{then}
                    $\hat{\nu} \gets 4 + 6/(\hat{\kappa} - 3)$
                    \textbf{else} $\hat{\nu} \gets 100$
\STATE \quad\quad $\nu \gets \eta\,\nu + (1-\eta)\,\text{clip}(\hat{\nu},\, 3,\, 100)$

\STATE \quad \textbf{(S2) Draw latent precision weights $\boldsymbol{\lambda}$}
\STATE \quad\quad \textbf{if} $t = 1$ \textbf{then}
\STATE \quad\quad\quad $\boldsymbol{\lambda} \gets \text{Gamma}(\text{shape}=\nu/2,\,
                         \text{rate}=\nu/2)$
                         \hfill $\triangleright$ prior
\STATE \quad\quad \textbf{else}
\STATE \quad\quad\quad $\mathbf{r}^2 \gets \hat{\mathbf{r}}^2$
\STATE \quad\quad\quad $\boldsymbol{\lambda} \gets \text{Gamma}(\text{shape}=(\nu+1)/2,\,
                         \text{rate}=(\nu+\mathbf{r}^2)/2)$
                         \hfill $\triangleright$ posterior

\STATE \quad \textbf{(S3) Compute perturbation scaling}
\STATE \quad\quad $\mathbf{w} \gets 1 \;/\; \sqrt{\max(\boldsymbol{\lambda},\, 10^{-6})}$

\STATE \quad \textbf{Perturb observations}
\STATE \quad\quad \textbf{for} $j = 1$ \textbf{to} $N_e$ \textbf{do}
\STATE \quad\quad\quad $\tilde{\mathbf{D}}_{:,j} \gets \mathbf{y} +
                         \sqrt{\alpha_t}\; \boldsymbol{\phi}_{:,j} \odot \mathbf{w} \odot
                         \mathcal{N}(\mathbf{0}, \mathbf{1})$

\STATE \quad \textbf{Kalman update}
\STATE \quad\quad $\Delta\mathbf{D} \gets \mathbf{D} - \text{mean}(\mathbf{D}, \text{axis}=1)$
\STATE \quad\quad $\Delta\mathbf{M} \gets \mathbf{M} - \text{mean}(\mathbf{M}, \text{axis}=1)$
\STATE \quad\quad $\mathbf{C}_{md} \gets \Delta\mathbf{M} \,
                          \Delta\mathbf{D}^{\mathsf{T}} \,/\, (N_e - 1)$
\STATE \quad\quad $\mathbf{C}_{dd} \gets \Delta\mathbf{D} \,
                          \Delta\mathbf{D}^{\mathsf{T}} \,/\, (N_e - 1)$
\STATE \quad\quad \textbf{for} $j = 1$ \textbf{to} $N_e$ \textbf{do}
\STATE \quad\quad\quad $\mathbf{C}_d \gets \text{diag}(\boldsymbol{\phi}_{:,j}^2 \odot
                             \mathbf{w}^2)$
                             \hfill $\triangleright$ inflated by $1/\boldsymbol{\lambda}$
\STATE \quad\quad\quad $\mathbf{K} \gets \mathbf{C}_{dd} + \alpha_t \mathbf{C}_d$
\STATE \quad\quad\quad $\mathbf{M}_{:,j} \gets \mathbf{M}_{:,j} +
                             \mathbf{C}_{md} \, \mathbf{K}^{-1} \,
                             (\tilde{\mathbf{D}}_{:,j} - \mathbf{D}_{:,j})$

\STATE \quad \textbf{Update $\boldsymbol{\phi}$ (IKEA Laplace)}
\STATE \quad\quad $\boldsymbol{\phi} \gets \text{covariance\_matrix}(
                         \mathbf{D},\; \text{obs\_data},\; \alpha_t,\; \ldots)$

\STATE \quad \textbf{Report:}
         $\text{mean}(\boldsymbol{\phi})$, $\nu$,
         $\text{mean}(\tilde{\mathbf{M}}, \text{axis}=1)$

\STATE \textbf{end for}
\STATE \textbf{return} $\mathbf{M}$, $\nu$, $\boldsymbol{\phi}$
\end{algorithmic}
\end{algorithm}

%=============================================================================
\section{Implementation Notes}
%=============================================================================

\subsection{Inversion Type Compatibility}

The perturbation and covariance scaling (S3) must be applied to all
inversion types supported by the base ES-MDA implementation:
\texttt{efast\_subspace}, \texttt{fast\_subspace}, \texttt{subspace},
\texttt{svd}, \texttt{esmda}, \texttt{esmda\_dask}, and \texttt{dask}.
For the subspace-family methods, $\mathbf{w}^2$ enters the per-member
noise vector $\mathbf{a}_{C_d} = (N_e-1)\,\alpha_t\,
\boldsymbol{\phi}_{:,j}^2 \odot \mathbf{w}^2$.
For the eFAST and Dask methods, the condensed noise vector
$\boldsymbol{\phi}_{\text{rand}}$ is multiplied element-wise by
$\mathbf{w}$.
For the SVD method, the diagonal of $\mathbf{C}_d$ is multiplied by
$\mathbf{w}^2$.

\subsection{Numerical Guards}

\begin{itemize}
  \item $\boldsymbol{\lambda}$ is clipped to $[10^{-6}, \infty)$ before
        computing $\mathbf{w}$ to prevent division by zero.
  \item Residuals are clipped to $\pm 8$ before kurtosis computation
        to robustify the $\nu$ estimator against individual extreme
        outliers.
  \item $\nu$ is clamped to $[3, 100]$: the lower bound ensures the
        Gamma shape parameter $\nu/2$ remains $\geq 1.5$, preventing
        degenerate draws; the upper bound provides a smooth transition
        to the Gaussian limit.
  \item The EMA smoothing constant $\eta = 0.7$ was chosen heuristically
        and has not been systematically optimised. Lower values produce
        faster adaptation but more iteration-to-iteration variation in
        $\nu$.
\end{itemize}

\subsection{Computational Cost}

The Student-$t$ extension adds three vector operations per iteration:
a Gamma draw of length $N_d$ ($\sim$2~ms for $N_d = 11{,}323$), a
kurtosis computation (one pass over $N_d$ residuals), and an EMA update
(scalar).
Total overhead is under 0.1~seconds per iteration, compared to 8--12
seconds for the forward model evaluations and Kalman update.
The asymptotic complexity of the ES-MDA loop is unchanged.

\subsection{Relationship to the Gaussian Limit}

As $\nu \to \infty$, the Gamma distribution concentrates at its mean,
$\lambda_i \to 1$ for all $i$, and $\mathbf{w} \to \mathbf{1}$.
The algorithm reduces exactly to the standard Gaussian ES-MDA.
This property enables a single code path for both likelihoods, with
the Gaussian case recovered by setting $\mathbf{w} = \mathbf{1}$
(equivalent to \texttt{likelihood="gaussian"} in the implementation).

\begin{thebibliography}{99}

\bibitem[Botha et al.(2023)]{Botha2023}
  Botha~I, Adams~MP, Tran~DH, Bennett~FR, Drovandi~CC (2023).
  \newblock ``Component-Wise Iterative Ensemble Kalman Inversion for
             Static Bayesian Models.''
  \newblock \emph{Statistics and Computing}, \textbf{33}, 109.

\bibitem[Emerick and Reynolds(2013)]{Emerick2013}
  Emerick~AA, Reynolds~AC (2013).
  \newblock ``Ensemble Smoother with Multiple Data Assimilation.''
  \newblock \emph{Computers \& Geosciences}, \textbf{55}, 3--15.

\bibitem[West(1987)]{West1987}
  West~M (1987).
  \newblock ``On Scale Mixtures of Normal Distributions.''
  \newblock \emph{Biometrika}, \textbf{74}(3), 646--648.

\end{thebibliography}

\end{document}
